\section{UCRA: Methodology}
\label{sec:method}

UCRA organizes reservation control into five stages wrapped in one
closed loop (Fig.~\ref{fig:pipeline}): uncertainty \emph{estimation}
(Stage~1), uncertainty-to-reservation \emph{transformation}
(Stage~2), risk-adaptive \emph{allocation} (Stage~3), reservation
\emph{update and release} (Stage~4), and \emph{self-evolving
learning} (Stage~5). The design principle throughout is that learned
components forecast, while committed decisions remain auditable
arithmetic. We detail each stage below; Table~\ref{tab:notation}
summarizes the notation.

\begin{table}[!t]
\caption{Principal notation.}
\label{tab:notation}
\centering
\small
\begin{tabular}{@{}p{0.16\linewidth}p{0.74\linewidth}@{}}
\toprule
Symbol & Meaning \\
\midrule
$\delta$ & Slot duration (15 min) \\
$K$ & Input window length (96 slots = 24 h) \\
$H$ & Forecast horizon (4 slots = 1 h) \\
$\D_t$ & Realized aggregate demand in slot $t$ \\
$C$ & Physical capacity of the pool \\
$\qtau$ & Forecast $\tau$-quantile of next-slot demand \\
$\uhat$ & Point forecast (median quantile) \\
$\spread$ & Forecast dispersion $\hat q_{t,0.99} - \hat q_{t,0.50}$ \\
$\rhot$ & Online risk indicator in $[0, 1.5]$ \\
$\kappa$ & Risk dial: buffer size in units of spread \\
$\Rt$ & Committed reservation for slot $t$ \\
$V, U, W, E$ & Violation, utilization, waste, tracking error \\
$\mathcal{B}$ & Reservoir replay buffer (capacity 512) \\
\bottomrule
\end{tabular}
\end{table}

\begin{figure}[!t]
\centering
\diagramplaceholder{5.2cm}{ucra-pipeline}
{Suggested content: a horizontal five-stage data-flow: \{24 h demand
window\} $\rightarrow$ \textbf{S1 Quantile LSTM} $\rightarrow$
$(\qtau, \spread, \rhot)$ $\rightarrow$ \textbf{S2 $\Phi$ transform}
$\rightarrow$ $R^{target}_t$ $\rightarrow$ \textbf{S3 Allocation}
$\rightarrow$ \textbf{S4 EWMA + hysteresis} $\rightarrow$ \textbf{network}
$\rightarrow$ feedback ($\D_t$, violation flags) splitting into
$\rhot$-state (to S2) and \textbf{S5 drift monitor $\rightarrow$
DelayedReplay $\rightarrow$ fine-tune} (back to S1).}
\caption{The UCRA closed loop. Feedback from committed reservations
feeds both the risk indicator (Stage~2) and the drift-triggered
self-evolution path (Stage~5).}
\label{fig:pipeline}
\end{figure}

\subsection{Stage 1: Quantile Demand Forecasting}
\label{sec:stage1}

The forecasting problem is: given the demand window of the last $K$
slots, output the predictive distribution of demand for each of the
next $H$ slots. We use a two-layer LSTM with 64 hidden units and
dropout 0.1~\cite{hochreiter1997lstm}, reading the raw demand series
one value per timestep, followed by a linear head that emits
$H \times Q$ values --- one quantile per level
$\tau \in \{0.05, 0.25, 0.50, 0.75, 0.90, 0.95, 0.99\}$ per horizon
step --- totaling 52{,}252 parameters. The network is deliberately
standard: forecast \emph{accuracy} is not where UCRA's contribution
lies, and a widely understood architecture keeps the audited
information flows simple.

Training minimizes the multi-quantile pinball
loss~\cite{koenker1978regression}. For predictions
$\hat q^{(\tau)}_h$ and realized demand $y_h$ at horizon step $h$:
\begin{equation}
\mathcal{L}_{\text{pin}} \;=\; \frac{1}{H}\sum_{h=1}^{H}
\frac{1}{Q}\sum_{j=1}^{Q}
\rho_{\tau_j}\!\bigl(y_h - \hat q^{(\tau_j)}_h\bigr),
\qquad
\rho_\tau(e) = \max\bigl(\tau e,\, (\tau-1)e\bigr).
\label{eq:pinball}
\end{equation}
Both inputs and targets are $z$-scored with statistics from the
\emph{training segment only}; this matters because pinball gradients
are bounded by $\tau$ and would stall against targets in the tens of
thousands. Training uses Adam (learning rate $10^{-3}$, batch 64)
with early stopping on validation pinball loss (patience 6, maximum
50 epochs) on a chronological 70/15/15 split --- no shuffling across
time. At inference the scaler is inverted so that all downstream
stages operate in original demand units.

From the per-slot forecast vector, Stage~1 assembles the triple that
the rest of the loop consumes:
\begin{align}
\uhat &= \hat q_{t,0.50}^{(1)} \quad \text{(next-slot median)}, \label{eq:uhat}\\
\spread &= \max\bigl(0,\; \hat q_{t,0.99}^{(1)} - \hat q_{t,0.50}^{(1)}\bigr),
\quad \text{(dispersion)}, \label{eq:spread}\\
\rhot &= \min\!\Bigl(1.5,\; \bar v_t + \tfrac{1}{4}\max\bigl(0,\, z_t - 2\bigr)\Bigr),
\label{eq:rho}
\end{align}
where $\bar v_t$ is the rolling violation rate over the last 96 slots
and $z_t = |\bar\D_t - \mu_{\text{train}}| / \sigma_{\text{train}}$
the deviation of the recent mean demand from training statistics.
$\rhot$ is thus a bounded, two-channel risk indicator: it rises with
\emph{operational} distress (violations) and with \emph{distributional}
distress (demand drifting away from the regime the forecaster was
trained on). Note that $\rhot$ is computed exclusively from slots
$< t$ --- a first instance of the information-flow discipline that
Section~\ref{sec:stage5} generalizes.

\subsection{Stage 2: The Uncertainty-to-Reservation Transform
$\Phi$}
\label{sec:stage2}

The transform maps the forecast triple to a reservation target. The
base term follows the forecast's own $\tau$-quantile, interpolated
from the median and the dispersion,
\begin{equation}
\qtau \;=\; \uhat + \frac{\tau - 0.5}{0.99 - 0.5}\,\spread,
\qquad \tau = 0.9,
\label{eq:qtau}
\end{equation}
and the risk buffer scales with dispersion and risk feedback:
\begin{equation}
R^{\text{target}}_t \;=\; \operatorname{clip}\!\Bigl(
\qtau + \kappa\,\spread\,\bigl(1 + w\,\rhot\bigr),\;
(1 + h_{\min})\,\uhat,\;
\gamma_{\max} C\Bigr),
\label{eq:phi}
\end{equation}
with the operating point $\tau = 0.9$, $\kappa = 0.5$, $w = 0.3$,
headroom floor $h_{\min} = 0.05$, and ceiling
$\gamma_{\max} = 0.95$. Each term has a distinct role and is
individually ablatable. The base term $\qtau$ makes the reservation
follow demand. The buffer term $\kappa \spread$ prices forecast
\emph{uncertainty}: when the model is confident, the reservation
tightens toward the median; when the band widens (night-to-morning
transitions, unstable regimes), the buffer widens in proportion. The
modulation $(1 + w \rhot)$ is the feedback term: under sustained
violations or demand drift the entire buffer inflates by up to
$1 + 0.3 \times 1.5 = 45\%$, and relaxes back when risk recedes. The
clamps encode two engineering floors: never reserve less than 5\%
above the point forecast (protection against pathologically confident
forecasts), and never reserve more than 95\% of capacity (the
remaining 5\% is control headroom).

$\kappa$ is the operator's risk dial. Its semantics are exact: an
additional unit of $\kappa$ reserves one additional $\spread$ of
buffer, everywhere, always. Section~\ref{sec:results-kappa} traces
the induced frontier from $\kappa = 0$ (bare quantile tracking) to
$\kappa = 2$; the knee --- the smallest zero-violation operating
point --- lies between 0 and 0.25 on our data, and $\kappa = 0.5$
provides margin at a measured utilization cost.

\subsubsection{Why calibrated quantiles translate into violation
control}

The link between Stage-1 calibration and the reservation's violation
rate can be made explicit, and it explains why the deployed operating
point achieves zero violations while a raw historical quantile does
not.

\begin{proposition}[Reservation violation under calibration]
\label{prop:bound}
Let $D_{t+1}$ be the next-slot demand and let the Stage-1 forecast be
perfectly calibrated in the quantile sense: for every level $u$,
$\Pr\bigl(D_{t+1} \le \hat q_{t,u} \mid \mathcal{F}_t\bigr) = u$.
If the reservation is set to
$R = \hat q_{t,\tau} + \kappa\,\spread\,(1 + w\,\rhot)$ with buffer
term $B = \kappa\,\spread\,(1 + w\,\rhot) \ge 0$ deterministic given
$\mathcal{F}_t$, then the conditional violation probability is
$\Pr\bigl(D_{t+1} > R \mid \mathcal{F}_t\bigr) \le 1 - \tau$.
\end{proposition}
\begin{IEEEproof}
Since $B \ge 0$ is $\mathcal{F}_t$-measurable, the event
$\{D_{t+1} > \hat q_{t,\tau} + B\}$ is a subset of
$\{D_{t+1} > \hat q_{t,\tau}\}$, whose conditional probability is, by
calibration, exactly $1 - \tau$. Monotonicity of probability yields
the claim; equality holds when $B = 0$.
\end{IEEEproof}

Two practical readings follow. First, calibration alone bounds the
violation rate by the tail mass above $\tau$ (10\% at $\tau = 0.9$);
pushing the realized rate to zero --- as our deployment point does
--- requires buffer beyond $\tau$, which is exactly what the kappa
term buys and what the frontier of Section~\ref{sec:results-kappa}
prices. Second, the bound degrades gracefully with miscalibration:
if the realized coverage at level $u$ is $\tilde u$ (e.g., our
q05--q99 coverage of 93.7\% against nominal 94\%), the violation
bound tracks $1 - \tilde\tau$ with the seed-level spread we measure
in Section~\ref{sec:results-seeds}. The comparison policies fail
precisely at these two points: the median forecast sets $\tau =
0.5$ (violation rate $1 - 0.5 = 50\%$, realized 48.0\%), and the
historical rolling quantile is not a calibrated conditional quantile
of $D_{t+1}$ at all (realized 16.3\% even for the test-window
oracle).

\subsubsection{Worked example of the transform}

For intuition, one concrete slot from the test window. Suppose the
forecast returns $\uhat = 61{,}200$ (median), $\spread = 12{,}400$
(q99 $-$ q50), and the risk state is calm, $\rhot = 0$. The base
term is $\qtau = 61{,}200 + \frac{0.4}{0.49} \cdot 12{,}400 =
71{,}322$; the buffer adds $\kappa \spread = 6{,}200$, giving
$77{,}522$ before clamping; the floor $(1.05 \times 61{,}200 =
64{,}260)$ does not bind, the ceiling $0.95 \times 143{,}954 =
136{,}757$ does not bind; Stage~4 then EWMA-blends toward this
target. In the same slot under a stressed risk state
($\rhot = 1.5$), the buffer inflates by a factor $1 + 0.3 \times
1.5 = 1.45$, reserving $6{,}200 \rightarrow 8{,}990$ extra ---
a targeted, reversible, and fully explainable adjustment. Every
reservation decision decomposes into these four readable terms.

\subsection{Stage 3: Risk-Adaptive Allocation}
\label{sec:stage3}

When the reservable pool is shared by $n$ entities (sectors or
slices) with current demand estimates $d_1, \dots, d_n$ and risks
$r_1, \dots, r_n$, the aggregate reservation $\Rt$ is split by
pressure-and-risk weights,
\begin{equation}
\omega_i = d_i^{\,p} \,(1 + r_i)^{\,s},
\qquad
A_i = \min\!\Bigl(d_i,\; (1-\beta)\,\Rt\,
\frac{\omega_i}{\sum_j \omega_j}\Bigr),
\label{eq:alloc}
\end{equation}
with pressure exponent $p = 1$, risk exponent $s = 0.5$, and a
best-effort share $\beta = 0.1$ withheld from all guaranteed
allocations. Leftover capacity from entities whose demand falls below
their share is re-distributed proportionally to the still-unmet
entities; demand beyond an allocation counts as unmet. The design
intent is that URLLC-class entities (high $r_i$) systematically
receive a larger multiplicative share of any reservation than their
raw pressure alone would justify, while $\beta$ preserves a small
best-effort pool so that guarantees never consume 100\% of the
committed capacity.

\subsection{Stage 4: Reservation Update and Release}
\label{sec:stage4}

Committing $R^{\text{target}}_t$ directly would make the reservation
chase every forecast wiggle. Stage~4 smooths the commitment with an
EWMA,
\begin{equation}
R_{t+1} = (1-a)\,\Rt + a\,R^{\text{target}}_t,
\qquad a = 0.3,
\label{eq:ewma}
\end{equation}
and adds \emph{release hysteresis}: downward adjustments are
accelerated only when the previous reservation left substantial
unused headroom,
\begin{equation}
R_{t+1} = (1-a)\,\Rt + a\,\max\bigl(R^{\text{target}}_t,\;
\D_t\bigr)
\quad \text{if } \frac{\Rt - \D_t}{C} > \theta_{\text{rel}}
\text{ and } R^{\text{target}}_t < \Rt,
\label{eq:hysteresis}
\end{equation}
with release threshold $\theta_{\text{rel}} = 0.1$. The effect is
asymmetric by design: ramping up is responsive (violations are
immediate SLA damage), while ramping down requires sustained evidence
that the slack is truly unused, which suppresses the
expand--contract oscillation that bursty demand otherwise induces.
The loop starts from the training-window 90\% quantile, a
deliberately conservative cold start.

\subsection{Stage 5: Self-Evolving Learning Under a Causal
Information-Flow Contract}
\label{sec:stage5}

Stage~5 is UCRA's insurance policy. Its premise: Stages~1--4 assume
the forecaster's picture of demand matches reality; when the world
drifts --- a mass event, a new traffic mix --- the forecast silently
stales, $\Phi$ faithfully reserves around the \emph{wrong}
distribution, and violations climb. Stage~5 watches the feedback
stream, detects that climb, and repairs the forecaster --- while
never pretending to know things it cannot yet know.

\subsubsection{Drift monitoring}

A dual-trigger monitor runs over the test stream. Trigger~A is
operational: the rolling violation rate over the last 96 slots
exceeds $v^* = 0.15$. Trigger~B is distributional: the current demand
z-score against training statistics exceeds $z^* = 3.0$. Triggers are
evaluated every $E_{\text{ev}}$ slots (96 in the steady-state run,
24 in the surge experiment) and, by construction, only on information
already observed. The two channels are deliberately redundant: A
fires when the model is wrong \emph{in the way that hurts users}; B
fires when the world has moved even before the reservations feel it.

\subsubsection{Reservoir replay}

On trigger, the engine fine-tunes the forecaster on a replay buffer
$\mathcal{B}$ of (window, label) pairs maintained by reservoir
sampling~\cite{vitter1985random}: the first 512 pairs fill the
buffer; each subsequent pair replaces a uniformly chosen slot with
probability $512/n_{\text{seen}}$. At any moment, every past window
had an equal inclusion probability, so the buffer holds both
pre-drift and post-drift windows --- the standard replay defense
against catastrophic forgetting~\cite{chaudhry2019continual,
rolnick2019experience}. Fine-tuning runs 3 epochs of full-batch Adam
at learning rate $0.3 \times 10^{-3}$ (gentler than training) with
the same pinball objective as Stage~1.

\subsubsection{The causal contract}

\begin{proposition}[Causal label availability]
\label{prop:causal}
Let $H$ be the forecast horizon. Suppose a training window whose
labels cover slots $s \in \{w, \dots, w+H-1\}$ is offered to the
buffer at loop time $w$, and admitted to the replay pool only at the
first loop time $t$ satisfying $t \ge w + H - 1$. Then any update
executed at loop time $t$ trains exclusively on labels
$D_s$ with $s \le t$ --- labels that a live deployment would have
already observed.
\end{proposition}
\begin{IEEEproof}
The window offered at $w$ carries labels for slots $w$ through
$w + H - 1$; at time $w$ only label $D_w$ has been observed. Under
the admission rule, admission at $t$ requires $w + H - 1 \le t$,
hence every label slot $s \le w + H - 1 \le t$. Updates at time $t$
sample only admitted windows, so every label they touch satisfies
$s \le t$.
\end{IEEEproof}

Two implementation rules complete the contract. First,
\emph{delayed replay intake}: our pipeline realizes the admission
rule with a staging queue (\texttt{DelayedReplay}) between the live
loop and the reservoir buffer; an offered window waits there
$H - 1$ slots before becoming trainable. Second, \emph{partial
forecast refresh}: after an update at time $t$, only forecasts for
slots $t' > t$ may change. Slots $t' \le t$ were already served with
specific reservations; rewriting their forecast post hoc would (i)
corrupt any statistic computed from the served forecast and (ii)
misrepresent what the system actually did. Both rules are enforced
in code and covered by dedicated regression tests
(Section~\ref{sec:setup}); the leaky variants of both rules existed
in our own earlier implementation and are quantified as
\emph{L1} and \emph{L2} in Section~\ref{sec:leakage-audit}.

\begin{figure}[!t]
\centering
\diagramplaceholder{4.6cm}{causal-replay-timeline}
{Suggested content: a horizontal time axis of test slots $t = 0
\dots 15$ with $H = 4$. For one window offered at $t = 5$, draw its
four label slots $5..8$ and shade slots $6..8$ as ``not yet
observable'' until slot 8; mark ``admitted to replay'' at slot
$t + H - 1 = 8$; mark a trigger-fired fine-tune at slot 12 whose
training batch reaches only labels $\le 12$; and show the
forecast-refresh boundary at the update slot (slots $\le$ update
keep the served forecast).}
\caption{Causal replay timeline: a window offered at $t$ becomes
trainable exactly at $t + H - 1$, and an update at $t$ only moves
forecasts strictly beyond $t$.}
\label{fig:causal-timeline}
\end{figure}

\subsubsection{Complexity}

Per slot, the loop performs one LSTM forward pass
($O(K)$ in the window length, 96 steps), $O(1)$ transform
arithmetic, and $O(n)$ allocation --- microseconds in practice on
CPU. A fine-tune event processes $3 \times 512$ window-labels at
$O(K)$ each --- on the order of one training epoch over a small
dataset, triggered at most every $E_{\text{ev}}$ slots. The memory
footprint is the 512-window buffer ($\approx$2\,MB in float32).
Stage~5 is thus deployable on commodity edge hardware; no GPU is
required anywhere in the loop at inference time.
